% answers_short.tex —— 2012 年·陕西卷·理科 · 简短答案速查
% 由 tools/make_short_answers.py 自 answers.tex 的 \gksolhead/\gkanswer 提取，
% 并与 2012/2012_数学真题_陕西理_答案要点.md 逐题比对后生成，勿手改（改源 MD 后重跑）。

\begin{gknotice}
\textbf{使用说明}：本节只给最终答案便于快速核对，完整推导见解析版。答案为编写者独立推导，\textbf{非官方标准答案}。
\end{gknotice}

\gkpart{选择题}
\begin{center}\begin{tabular}{ccccccccccc}
\toprule
题号 & 1 & 2 & 3 & 4 & 5 & 6 & 7 & 8 & 9 & 10 \\
\midrule
答案 & C & D & B & A & A & B & D & C & C & D \\
\bottomrule
\end{tabular}\end{center}

\gkpart{填空题}
\begin{itemize}[leftmargin=2.6em]
  \item[\textbf{11.}] $1+\dfrac1{2^2}+\dfrac1{3^2}+\dfrac1{4^2}+\dfrac1{5^2}<\dfrac95$。
  \item[\textbf{12.}] $1$。
  \item[\textbf{13.}] $2\sqrt6$。
  \item[\textbf{14.}] $2$。
  \item[\textbf{15.}] A $[-2,4]$；B $5$；C $\sqrt3$。
\end{itemize}

\gkpart{解答题}
\begin{itemize}[leftmargin=2.6em]
  \item[\textbf{16.}] (1) $f(x)=2\sin\left(2x-\dfrac\pi6\right)+1$；(2) $\alpha=\dfrac\pi3$
  \item[\textbf{17.}] (1) $q=-2$；(2) 证明见解析
  \item[\textbf{18.}] (1) 证明见解析；(2) 逆命题为“若 $a\perp c$，则 $a\perp b$”，是真命题
  \item[\textbf{19.}] (1) $\dfrac{x^2}4+\dfrac{y^2}{16}=1$；(2) $y=x$ 或 $y=-x$
  \item[\textbf{20.}] (1) $0.22$；(2) $P(X=0)=0.5$，$P(X=1)=0.49$，$P(X=2)=0.01$；$EX=0.51$
  \item[\textbf{21.}] (1) 证明见解析；(2) $[-2,2]$；(3) $\{x_n\}$ 是递增数列
\end{itemize}
